\section{Arsitektur Buku dan Alur Belajar Tujuh Belas Bab}
\label{sec:ch01-arsitektur}

Buku ini disusun dalam tiga lapis. Lapis pertama adalah orientasi dan
landasan teori (Bab I--II). Lapis kedua adalah empat belas bab inti
(Bab III--XVI) yang mengikuti urutan topik RPS. Lapis ketiga adalah
evaluasi dan integrasi kompetensi (Bab XVII). Lampiran A--D memuat
rubrik, contoh tugas, glosarium, dan sumber belajar tambahan.

\begin{figure}[H]
\centering
\begin{tikzpicture}[node distance=4mm and 6mm, scale=0.82, transform shape]
  \node[phase, text width=2.1cm] (a) {Orientasi\\dan Landasan};
  \node[phase, text width=2.1cm, right=of a] (b) {Analisis\\Leksikal};
  \node[phase, text width=2.1cm, right=of b] (c) {Analisis\\Sintaksis};
  \node[phase, text width=2.1cm, right=of c] (d) {Analisis\\Semantik};
  \node[phase, text width=2.1cm, right=of d] (e) {Evaluasi\\dan Integrasi};
  \draw[arrow] (a) -- (b);
  \draw[arrow] (b) -- (c);
  \draw[arrow] (c) -- (d);
  \draw[arrow] (d) -- (e);
\end{tikzpicture}
\caption{Alur besar buku: dari orientasi menuju integrasi kompetensi}
\label{fig:ch01-alur-buku}
\end{figure}

Urutan bab dirancang agar setiap bab memakai hasil bab sebelumnya.
Sebagai contoh, Bab V memperkenalkan token, Bab VI memodelkan pengenal
token dengan otomata, dan Bab VII mengimplementasikan keduanya menjadi
\eng{lexer} yang bekerja. Pola yang sama berlanjut pada parsing,
analisis semantik, dan validasi AST.

Setiap bab inti memuat lima bagian materi, contoh berpenyelesaian,
praktikum, latihan, asesmen, daftar cek, rangkuman, dan evaluasi.
Komponen tersebut bukan pelengkap, melainkan bukti belajar yang dapat
dinilai.
